\section{Methodology}
\label{sec:method}

\subsection{Three-Tier Architecture Overview}

Figure~\ref{fig:architecture} illustrates the three-tier architecture. Each tier has distinct inputs, processes, and outputs:

\begin{figure}[t]
\centering
\begin{verbatim}
                    ┌───────────────┐
                    │  BOARD TIER   │  Monorepo level
                    │ (B1/B2/B3)    │  Cross-module synthesis
                    └───────┬───────┘
                        ↕ findings bubble up
                        ↕ revisions flow down
                    ┌───────────────┐
                    │  PANEL TIER   │  Module level
                    │ (PP1/PP2/PP3) │  Cross-portfolio panel
                    └───────┬───────┘
                        ↕ findings bubble up
                        ↕ revisions flow down
              ┌─────────────────────────┐
              │      PAPER TIER         │  Individual paper level
              │  (P1/P2/P3)             │  Per-paper reviews
              └─────────────────────────┘
\end{verbatim}
\caption{Three-tier hierarchical review architecture with bidirectional flow.}
\label{fig:architecture}
\end{figure}

\subsection{Paper Tier: Individual Review and Synthesis}

The paper tier follows the 8-stage lifecycle from prior work~\cite{dellalibera2026review}: draft → panel → synthesis → revision → recheck → ready → submit → accepted. For each paper:

\textbf{Inputs}: Paper LaTeX source, target venue, research domain.

\textbf{Process}:
\begin{enumerate}
\item Assemble reviewer panel (5-7 experts) matched to domain and venue~\cite{dellalibera2026calibration}.
\item Generate structured reviews (REVIEW-EXPERT-NAME.md).
\item Synthesize reviews into SYNTHESIS.md with P1/P2/P3 priority classification.
\item Generate REVISION-PLAN.md with specific edit instructions.
\item Apply revisions to LaTeX source~\cite{dellalibera2026dynamics}.
\item Iterate until convergence (avg score ≥ 2.5/4, no score < 2/4).
\end{enumerate}

\textbf{Outputs}:
\begin{itemize}
\item SYNTHESIS.md with P1/P2/P3 items
\item Revised paper LaTeX
\item Completion signal (ready for panel tier)
\end{itemize}

\textbf{Priority classification at paper tier}:
\begin{itemize}
\item \textbf{P1}: 3+ reviewers flag the issue OR issue threatens paper acceptance
\item \textbf{P2}: 2+ reviewers flag the issue
\item \textbf{P3}: 1 reviewer flags the issue
\end{itemize}

\subsection{Panel Tier: Cross-Portfolio Synthesis}

When all papers in a module reach the ready stage, the panel tier activates.

\textbf{Inputs}:
\begin{itemize}
\item All SYNTHESIS.md files from papers in the module
\item All P1/P2/P3 items across papers
\item Module research goals and constraints
\end{itemize}

\textbf{Process}:
\begin{enumerate}
\item Assemble 7-member expert panel (distinct from paper reviewers to avoid bias).
\item Each panel member reviews all papers in the module.
\item Panel discusses papers in a structured session, generating REVIEW\_PANEL.md.
\item Cross-paper pattern detection: identify issues appearing in 2+ papers.
\item Priority escalation: P1 items appearing in 3+ papers escalate to PP1.
\item Strategic assessment: identify module-level themes, gaps, and risks.
\end{enumerate}

\textbf{Outputs}:
\begin{itemize}
\item REVIEW\_PANEL.md with cross-portfolio synthesis
\item PP1/PP2/PP3 priorities per paper
\item Completion signal (module ready for board tier)
\end{itemize}

\textbf{Priority classification at panel tier}:
\begin{itemize}
\item \textbf{PP1}: Cross-paper pattern (3+ papers) OR issue threatens module coherence
\item \textbf{PP2}: Issue affects 2+ papers
\item \textbf{PP3}: Issue affects 1 paper but has module-level implications
\end{itemize}

\subsection{Board Tier: Program-Level Coordination}

When all modules complete panel reviews, the board tier activates.

\textbf{Inputs}:
\begin{itemize}
\item All REVIEW\_PANEL.md files from modules
\item All PP1/PP2/PP3 items across modules
\item Program research vision and strategic goals
\end{itemize}

\textbf{Process}:
\begin{enumerate}
\item Assemble 7-member board (program leadership, senior researchers).
\item Each board member reviews all modules.
\item Board meeting generates REVIEW\_BOARD.md with cross-module synthesis.
\item Cross-module pattern detection: identify issues appearing in 2+ modules.
\item Priority escalation: PP1 items affecting 3+ modules escalate to B1.
\item Strategic directive generation: board issues B1/B2/B3 directives per module.
\end{enumerate}

\textbf{Outputs}:
\begin{itemize}
\item REVIEW\_BOARD.md with program-level synthesis
\item B1/B2/B3 directives per module
\item Per-module BOARD-REVISION-PLAN-\{module\}.md files
\end{itemize}

\textbf{Priority classification at board tier}:
\begin{itemize}
\item \textbf{B1}: 3+ board members flag issue OR issue threatens program success
\item \textbf{B2}: Issue affects 2+ modules
\item \textbf{B3}: Issue affects 1 module but has program-level implications
\end{itemize}

\subsection{Bidirectional Flow Protocols}

\subsubsection{Upward Escalation (Issues Bubble Up)}

A paper-tier issue escalates to panel tier when:
\begin{enumerate}
\item Same issue appears in 3+ papers (P1 → PP1).
\item Issue appears in 2+ papers and threatens module goals (P2 → PP1).
\item Issue reveals systematic methodological weakness (any P-tier → PP1).
\end{enumerate}

A panel-tier issue escalates to board tier when:
\begin{enumerate}
\item Same issue appears in 3+ modules (PP1 → B1).
\item Issue affects 2+ modules and threatens program vision (PP2 → B1).
\item Issue reveals architectural inconsistency (any PP-tier → B1).
\end{enumerate}

\subsubsection{Downward Propagation (Directives Cascade)}

Board-level directives (B1/B2/B3) cascade as follows:
\begin{enumerate}
\item Board generates B1 directive (e.g., "Adopt consistent evaluation framework").
\item For each affected module, board generates BOARD-REVISION-PLAN-\{module\}.md with specific PP1/PP2 tasks.
\item Panel tier reviews plan, generates per-paper P1/P2 items in PANEL-REVISION-PLAN-\{paper\}.md.
\item Paper tier implements P1 items via LaTeX edits.
\end{enumerate}

This cascade ensures strategic decisions propagate to concrete code changes, closing the loop from vision to implementation.

\subsection{Round Management}

Each tier supports multiple review rounds:
\begin{itemize}
\item \textbf{Paper tier}: ROUND2-REVIEW-*.md, ROUND3-REVIEW-*.md, etc.
\item \textbf{Panel tier}: panel-reviews/round-1/, panel-reviews/round-2/, etc.
\item \textbf{Board tier}: board-reviews/round-1/, board-reviews/round-2/, etc.
\end{itemize}

Rounds at higher tiers trigger rounds at lower tiers: if board issues a B1 directive in round 2, affected modules enter panel round 2, which triggers paper round 3 for affected papers.

\subsection{Implementation}

The architecture is implemented as a Claude Code plugin with 11 commands:
\begin{itemize}
\item \texttt{panel:review} — Paper-tier lifecycle (8 stages)
\item \texttt{panel:convene} — Panel-tier cross-portfolio review
\item \texttt{panel:board} — Board-tier program coordination
\item \texttt{panel:status} — Cross-tier status dashboard
\item \texttt{panel:show} — Detailed paper/module/board view
\end{itemize}

The system uses YAML state files (\_panel.yaml per paper) to track stage, round, reviewers, scores, and priority items across all tiers.

\subsection{Review Aggregation Algorithm}
\label{sec:aggregation}

The synthesis process consolidates individual reviews into prioritized issue lists (P1/P2/P3, PP1/PP2/PP3, B1/B2/B3). This section details the aggregation algorithm addressing concerns raised in prior work~\cite{bernstein2015crowdsourcing,kamar2016complementarity}.

\subsubsection{Paper-Tier Aggregation (P1/P2/P3 Classification)}

Given $n$ reviews for a paper, where each review $R_i$ contains:
\begin{itemize}
\item Overall score $s_i \in \{1, 2, 3, 4\}$ (reject, weak accept, accept, strong accept)
\item Major issues $M_i = \{m_{i,1}, m_{i,2}, \ldots\}$ (blocking concerns)
\item Minor issues $m_i = \{n_{i,1}, n_{i,2}, \ldots\}$ (non-blocking concerns)
\end{itemize}

We extract and deduplicate issues across all reviews using semantic similarity clustering (embed issues with \texttt{text-embedding-3-large}, cluster at cosine similarity $\geq 0.80$). For each unique issue cluster $I$:

\textbf{P1 classification} (blocking):
\begin{equation}
\text{P1}(I) = \begin{cases}
\text{true} & \text{if } |R_i : I \in M_i| \geq 3 \\
\text{true} & \text{if } I \in M_i \text{ for any } i \text{ with } s_i = 1 \\
\text{true} & \text{if } \text{confidence}(I) > 0.85 \text{ and avg}(s_i : I \in M_i \cup m_i) < 2.0 \\
\text{false} & \text{otherwise}
\end{cases}
\end{equation}

\textbf{P2 classification} (important):
\begin{equation}
\text{P2}(I) = \neg\text{P1}(I) \land \left( |R_i : I \in M_i \cup m_i| \geq 2 \right)
\end{equation}

\textbf{P3 classification} (nice-to-have):
\begin{equation}
\text{P3}(I) = \neg\text{P1}(I) \land \neg\text{P2}(I) \land \left( |R_i : I \in M_i \cup m_i| = 1 \right)
\end{equation}

The confidence score $\text{confidence}(I)$ is computed as:
\begin{equation}
\text{confidence}(I) = \frac{|R_i : I \in M_i \cup m_i|}{n} \times \text{agreement}(I)
\end{equation}

where $\text{agreement}(I)$ is the mean pairwise cosine similarity of issue descriptions across reviews mentioning $I$. Low confidence ($< 0.60$) triggers human review (see Section~\ref{sec:oversight}).

\subsubsection{Panel-Tier Aggregation (PP1/PP2/PP3 Classification)}

At the panel tier, we aggregate across $k$ papers in a module. Let $P_j$ be paper $j$'s priority items. We identify cross-paper patterns using:

\textbf{Pattern detection}: Cluster all P1/P2/P3 items across papers (cosine similarity $\geq 0.75$). For each pattern $\Pi$ (cluster of items):

\textbf{PP1 classification} (module-critical):
\begin{equation}
\text{PP1}(\Pi) = \begin{cases}
\text{true} & \text{if } |\{j : \exists I \in P_j, I \in \Pi\}| \geq 3 \\
\text{true} & \text{if } \text{threatens\_module}(\Pi) = \text{true} \\
\text{true} & \text{if } \sum_{j} \text{P1}(I \in \Pi \cap P_j) \geq 5 \\
\text{false} & \text{otherwise}
\end{cases}
\end{equation}

where $\text{threatens\_module}(\Pi)$ is determined by panel review session (qualitative judgment).

\textbf{PP2 and PP3} follow analogous rules (2+ papers affected for PP2, 1 paper with module implications for PP3).

\subsubsection{Board-Tier Aggregation (B1/B2/B3 Classification)}

Board-tier aggregation operates on $m$ modules. Cross-module patterns are detected by clustering all PP1/PP2/PP3 items (similarity $\geq 0.70$). B1/B2/B3 classification follows the same structure as panel tier, with "modules" replacing "papers."

\subsubsection{Worked Example: P1 → PP1 Escalation}

Consider 5 papers in the merit module:
\begin{itemize}
\item Paper 1 (P2): "Evaluation lacks statistical power analysis"
\item Paper 2 (P2): "Results need significance testing"
\item Paper 3 (P1): "Missing statistical validation of claims"
\item Paper 4 (P3): "No confidence intervals reported"
\item Paper 5 (P1): "Statistical methods inadequately described"
\end{itemize}

These cluster at cosine similarity 0.82 into pattern $\Pi = $ "Statistical rigor issues." Since $|\{1,2,3,4,5\}| = 5 \geq 3$ and $\sum \text{P1}(\Pi) = 2$, this escalates to PP1: "Merit module papers lack consistent statistical validation framework."

\subsection{Human Oversight Mechanisms}
\label{sec:oversight}

To address concerns about AI-generated priority classifications amplifying errors~\cite{shneiderman2022human}, we implement three oversight layers:

\subsubsection{Confidence-Based Deferral}

All priority classifications include confidence scores (Equation~4). Classifications with $\text{confidence}(I) < 0.60$ are escalated to human review before affecting revision plans. The human reviewer:
\begin{enumerate}
\item Views the issue description and all source reviews
\item Sees the confidence score and contributing factors
\item Accepts, modifies, or rejects the classification
\item Optionally provides feedback to calibrate future classifications
\end{enumerate}

Empirically, 12\% of paper-tier issues and 8\% of panel-tier issues trigger deferral (see Section~\ref{sec:deferral-results}).

\subsubsection{Human Override Capability}

At each tier, humans can override priority classifications:
\begin{itemize}
\item \textbf{Paper tier}: Authors can contest P1 items via \texttt{\_panel.yaml} flag \texttt{p1\_override: "justification"}. A panel chair reviews overrides before the paper advances.
\item \textbf{Panel tier}: Module lead reviews all PP1 items before they propagate to papers, with authority to downgrade or defer.
\item \textbf{Board tier}: Board members vote on B1 items; unanimous consent required for B1 items affecting 3+ modules.
\end{itemize}

Override rates: 3\% at paper tier, 5\% at panel tier, 8\% at board tier (see Section~\ref{sec:override-results}).

\subsubsection{Emergent Pattern Validation}

For issues classified as "emergent" (not flagged at lower tiers), the system:
\begin{enumerate}
\item Surfaces the pattern with evidence (which papers/modules contribute)
\item Flags the pattern as "needs validation" if confidence $< 0.70$
\item Requires human approval before cascading to lower tiers as PP1/B1
\end{enumerate}

This prevents spurious emergent patterns (AI hallucinations) from triggering cascaded revisions.

\subsection{System Observability and Testing}
\label{sec:observability}

Production ML systems require observability, testing, and debugging infrastructure~\cite{shankar2022operationalizing}. We implement:

\subsubsection{Observability Metrics}

The system tracks 15 metrics across tiers:
\begin{itemize}
\item \textbf{Review quality}: Score distribution, consensus (std dev), outlier detection
\item \textbf{Priority classification}: P1/P2/P3 distribution per paper, escalation rates
\item \textbf{Synthesis quality}: Issue deduplication rate, cross-paper pattern frequency
\item \textbf{Revision efficiency}: P1 resolution rate per round, cascade completion time
\item \textbf{Confidence distribution}: Percentage of issues triggering deferral
\end{itemize}

Metrics are logged to \texttt{metrics.jsonl} per paper/module/board and aggregated into dashboards (see Figure~\ref{fig:dashboard} in Appendix). Anomaly detection flags:
\begin{itemize}
\item Reviewer scoring $> 2\sigma$ from mean (outlier)
\item Synthesis generating $> 15$ P1 items (excessive blocking issues)
\item Confidence scores trending downward over rounds (drift)
\end{itemize}

\subsubsection{Testing Strategy}

We maintain three test suites:
\begin{enumerate}
\item \textbf{Unit tests} (42 tests): Each stage handler (draft→panel, panel→synthesis, etc.) has tests verifying correct state transitions, priority classification, and error handling.
\item \textbf{Integration tests} (12 tests): Full lifecycle tests (draft→accepted) on synthetic papers with known-good expected outputs.
\item \textbf{Regression tests} (8 golden papers): Reference papers with manually validated reviews, used to detect when prompt/reviewer changes break existing functionality.
\end{enumerate}

Test coverage: 87\% of stage-machine code, 92\% of synthesis-engine code (measured via branch coverage).

\subsubsection{Debugging and Provenance}

Each priority item includes provenance metadata:
\begin{itemize}
\item \textbf{Source reviews}: Which reviewers flagged the issue
\item \textbf{Confidence breakdown}: How confidence was computed
\item \textbf{Escalation path}: If PP1/B1, which P/PP items contributed
\item \textbf{Trace ID}: Unique ID linking item through all tiers
\end{itemize}

The \texttt{panel:inspect <paper>} command displays full provenance for debugging. Structured logs (JSON) capture all intermediate states, enabling replay of synthesis on the same reviews to verify reproducibility.